\chapter{Ejecución de algoritmos}

\section{Construcción del horario global}

El horario global se construye a partir de los datos cargados, el grafo de conflictos y el solucionador principal.

\begin{lstlisting}[language=TypeScript,caption={Flujo principal del horario global}]
const input = SchedulerDataLoader.build().getData();
const { graph } = new ConflictGraphBuilder(input).build();
const algoritmo = new DSatur<GrupoData>(graph, input.salones.map(s => s.id));
const horario = new ScheduleBuilder(graph, input.salones).buildHorario(algoritmo);
\end{lstlisting}

\section{Uso de constraints durante la asignación}

Antes de aceptar una franja o un bloque, el sistema consulta los verificadores de restricciones. Cada restricción concreta extiende una clase abstracta y devuelve si la asignación es válida.

\begin{lstlisting}[language=TypeScript,caption={Consulta de restricciones}]
const checker = new AllocatorConstraintChecker(dayConstraints, slotConstraints);
const dayResult = checker.checkDay(dia, materia, asignadas, contexto);
const slotResult = checker.checkSlot(franja, materia, asignadas, true);
\end{lstlisting}

\section{Planificación estudiantil}

Cuando llega una solicitud de estudiante, el controlador valida el body, carga el horario global y delega la construcción al servicio de estudiantes.

\begin{lstlisting}[language=TypeScript,caption={Controlador de horarios de estudiante}]
async createSchedule(req: Request, res: Response): Promise<void> {
  const body = req.body;
  if (!body.estudianteId || !Array.isArray(body.materias) || body.materias.length === 0) {
    res.status(400).json({success: false, error: "Se requiere estudianteId y materias[]."});
    return;
  }

  const input = SchedulerDataLoader.build().getData();
  const horario = await scheduleService.getHorarioGlobal();
  const resultado = await studentService.createSchedule(horario, input.salones, body);
  res.json({success: true, data: resultado});
}
\end{lstlisting}

\section{Recuperación de datos ya calculados}

El servicio de horario global mantiene caché en memoria y el servicio de estudiantes guarda el resultado por ID. Eso permite devolver el horario sin recalcular todo el grafo en cada solicitud.

\begin{lstlisting}[language=TypeScript,caption={Caché de horario global}]
async getHorarioGlobal(): Promise<HorarioSemanal> {
  if (this.horarioGlobal) return this.horarioGlobal;
  const input = SchedulerDataLoader.build().getData();
  const { graph } = new ConflictGraphBuilder(input).build();
  const algoritmo = new DSatur<GrupoData>(graph, input.salones.map(s => s.id));
  const horario = new ScheduleBuilder(graph, input.salones).buildHorario(algoritmo);
  this.horarioGlobal = horario;
  return horario;
}
\end{lstlisting}
